% !PRINTED 500
\weberpage{500}
%<<CONT
11. \gesperrt{dass, wenn der Fall der Imprimitivität eintritt, die Gruppe $P$ im Falle 9. die ganze symmetrische Gruppe, im Falle 10. die ganze alternirende Gruppe für jedes einzelne System der Imprimitivität enthält.}

Die wiederholte Zusammensetzung einer und derselben Permutation mit
sich selbst bezeichnet man wie Potenzen. Wenn
also $\pi$ irgend eine Permutation ist, so muss unter den auf einander
folgenden Potenzen
\[
\pi^0, \pi^1, \pi^2, \pi^3 \ldots,
\]
worin $\pi^0 = 1$ die identische Permutation bedeutet, nothwendig
dieselbe Permutation wiederkehren. Aus $\pi^u = \pi^{u+e}$ folgt aber
$\pi^e = 1$. Ist $e$ der kleinste positive Exponent, für den $\pi^e = 1$
wird, so sind die Permutationen
\weq{4}{1, \pi, \pi^2 \ldots \pi^{e-1}}
alle von einander verschieden, während sie sich bei der Fortsetzung $\pi^e, \pi^{e+1}, \ldots \pi^{2e-1} \ldots$ in derselben Reihenfolge periodisch
wiederholen. Die Reihe (4) heisst die \gesperrt{Periode der Permutation}
$\pi$ und $e$ wird der \gesperrt{Grad von} $\pi$ genannt. Die Periode
von $\pi$ bildet eine Gruppe, da durch Zusammensetzung zweier
ihrer Elemente ein Element derselben Periode entsteht.
Ist $\pi$ eine cyklische Permutation
\[
\pi = (0, 1, 2 \ldots n - 1),
\]
so ist bei geradem $n$
\[
\pi^2 = (0, 2, 4 \ldots n - 2) (1, 3 \ldots n - 1),
\]
und bei ungeradem $n$
\[
\pi^2 = (0, 2, 4 \ldots n - 1, 1, 3 \ldots n - 2),
\]
und der Grad von $\pi$ ist in beiden Fällen gleich $n$.
Ist $\pi$ in mehrere Cyklen zerlegt, so ist der Grad von $\pi$
gleich dem kleinsten gemeinschaftlichen Vielfachen der Gliederzahl der einzelnen Cyklen, z.\,B.
\[
\pi = (0, 1) (2, 3, 4), \; \pi^2 = (0) (1) (2, 4, 3), \; \pi^3 = (0, 1) (2) (3) (4),
\]
% [sic] Der Druck trennt die disjunkten Cyklen in $\pi^4$ und $\pi^5$ durch Kommata, anders als in $\pi^2$ und $\pi^3$ unmittelbar davor; die gedruckte Notation bleibt erhalten.
\[
\pi^4 = (0), (1), (2, 3, 4), \; \pi^5 = (0, 1), (2, 4, 3), \; \pi^6 = 1.
\]
Wenn $\pi$ in irgend einer Permutationsgruppe $P$ vorkommt,
so enthält $P$ die ganze Periode von $\pi$.
